\chapter{Desenvolvimento Iterativo por Sprints e Validação por Testes Automatizados}

\section{Metodologia Ágil: Desenvolvimento por Sprints}

O desenvolvimento do sistema Python/Streamlit seguiu uma abordagem iterativa e
incremental, organizada em quatro \textit{sprints} sucessivas, alinhada aos
princípios do \textit{framework} Scrum~\cite{scrum2026}. Cada \textit{sprint}
produziu um incremento funcional e testável do sistema, permitindo validação
contínua da equivalência com o sistema Java legado antes do avanço para a
etapa seguinte.

\begin{description}

\item[a)] \textbf{\textit{Sprint} 1 --- Motor de Dados (\textit{Core Backend})}:
Esta etapa implementou as funções \texttt{parse\_matrix()} e
\texttt{parse\_vector()}, responsáveis pela leitura unificada de arquivos
\texttt{.txt}, \texttt{.csv} e \texttt{.xlsx}, com detecção automática de
separador de coluna e normalização do decimal brasileiro. A validação ocorreu
por meio de testes no console, sem envolvimento da interface Streamlit,
utilizando os \textit{datasets} reais \textit{EthanolLevel} e
\texttt{CGL\_nir}, confirmando as dimensões esperadas das matrizes.

\item[b)] \textbf{\textit{Sprint} 2 --- Algoritmos Científicos}: O núcleo
matemático foi codificado em \texttt{models.py} como funções puras,
independentes de qualquer biblioteca de interface. Implementaram-se
\texttt{perform\_rlm()} via \texttt{np.linalg.lstsq}, \texttt{perform\_bagging()}
para \textit{bagging} e \textit{subagging} parametrizados, \texttt{calc\_metrics()}
para RMSE, MAE e $R^2$, e \texttt{run\_cross\_validation()} para validação
cruzada 5-\textit{fold}, culminando na função \texttt{find\_optimal\_parameters()}
de \textit{grid search}.

\item[c)] \textbf{\textit{Sprint} 3 --- Interface Reativa}: Conectou-se
\texttt{models.py} ao orquestrador \texttt{app.py}, implementando os quatro
componentes \texttt{st.file\_uploader}, a validação de parâmetros por meio de
\texttt{parse\_int\_list()} e \texttt{parse\_float\_list()}, o \textit{callback}
de progresso e a persistência do resultado em \texttt{st.session\_state}. Os
quatro gráficos Plotly foram centralizados em \texttt{components.py}, garantindo
que nenhuma lógica de negócio residisse na camada de apresentação.

\item[d)] \textbf{\textit{Sprint} 4 --- Otimização com \textit{Cache}}: Avaliou-se
a arquitetura resultante quanto à sua compatibilidade com os mecanismos de
\textit{cache} do Streamlit. Como \texttt{parse\_matrix()}, \texttt{parse\_vector()}
e \texttt{run\_full\_analysis()} são funções puras, determinísticas e com
entradas serializáveis, o decorador \texttt{@st.cache\_data} pode ser aplicado
diretamente, sem refatoração, reduzindo a latência em execuções repetidas.

\end{description}

% ─────────────────────────────────────────────────────────────────────────────
\section{Suíte de Testes Automatizados com pytest}

Conforme indicado no TCC~I, a suíte de testes automatizados constitui a principal
entrega de validação deste TCC~II, fechando o ciclo entre a consolidação da base
algorítmica em Java --- realizada na etapa anterior --- e a comprovação formal de
sua equivalência comportamental na reescrita em Python/Streamlit apresentada a
seguir.

A correção do sistema Python foi validada por meio de uma suíte de 92 testes
automatizados com \texttt{pytest}~\cite{pytest2024}, organizada em três classes
complementares: testes unitários, testes de caixa-preta e testes de carga e
\textit{stress}. Essa estratificação permitiu verificar, respectivamente, a
corretude matemática isolada, a fidelidade em relação ao sistema Java legado e
o comportamento sob volumes de dados realistas.

\begin{description}

\item[a)] \textbf{Testes Unitários (21 testes)}: Exercitaram cada função pura
de \texttt{models.py} isoladamente, com matrizes de dimensão reduzida e valores
determinísticos, por meio das \textit{fixtures} \texttt{dados\_simples}
($20\times5$) e \texttt{dados\_medios}~($300\times50$). Cobriram-se oito testes
de \textit{parsing} de arquivos e treze testes de modelos de regressão,
incluindo a recuperação exata de $\boldsymbol{\beta}$ em sistemas determinados
e a reprodutibilidade do \textit{bagging} sob a mesma semente aleatória.

\item[b)] \textbf{Testes de Caixa-Preta (60 testes)}: Verificaram a equivalência
numérica entre o sistema Python e o sistema Java original, sem acesso ao
código-fonte legado, seguindo a estratégia de \textit{black-box
testing}~\cite{pressman2016}. Doze testes parametrizados foram aplicados a
cinco analitos (Mn, Mo, Cr, Ni, Fe), comparando RMSE, MAE e $R^2$ contra os
arquivos de referência gerados pelo sistema Java, com tolerância de
$5\times10^{-4}$ para o RLM determinístico e tolerâncias mais amplas para os
métodos estocásticos de \textit{bagging} e \textit{subagging}.

\item[c)] \textbf{Testes de Carga e \textit{Stress} (11 testes)}: Quantificaram
os limites operacionais de tempo de execução e consumo de memória para
matrizes espectrais de alta dimensionalidade~\cite{jorgensen2014}, empregando
as \textit{fixtures} \texttt{matriz\_grande}~($5.000\times723$) e
\texttt{matriz\_stress}~($15.000\times723$). O tempo foi medido com
\texttt{time.perf\_counter()} e o pico de alocação de memória com
\texttt{tracemalloc}, confirmando escalabilidade linear do \textit{bagging} em
função do número de \textit{bags}.

\item[d)] \textbf{Resultados Consolidados}: A suíte completa foi executada em
ambiente Windows com Python~3.12.10 e pytest~9.1.0, totalizando 92 testes
aprovados e nenhuma falha, em 81{,}63~segundos. Esse resultado evidencia três
conclusões objetivas: correção matemática comprovada nos testes unitários,
fidelidade numérica confirmada em relação ao sistema Java legado e capacidade
de carga validada para matrizes de até $15.000\times723$ sem erro de memória.

\end{description}
